\originalchapter{引力与天体轨道}\label{ch:gravity}
\originalsection{球壳、地球内部与测量}
均匀球壳半径 $a$, 壳面密度 $\sigma$, 观察点距中心 $r$. 沿观察轴取极角, 环形面元 $\dd m=2\pi\sigma a^2\sin\theta\dd\theta$, 到观察点距离 $s=\sqrt{r^2+a^2-2ar\cos\theta}$. 代 $s\dd s=ar\sin\theta\dd\theta$ 得势能
\begin{equation}
U(r)=-\frac{2\pi Gm\sigma a}{r}\int_{|r-a|}^{r+a}\dd s
=-\frac{GmM_s}{\max(r,a)}.
\end{equation}
壳外等效点质量, 壳内势能恒定、力为零. 把球体分成球壳, 得球外 $F=-GMm/r^2$, 球内仅内侧质量贡献, $F=-GmM(<r)/r^2$. 均匀密度 $\rho$ 时 $M(<r)=4\pi\rho r^3/3$, 因而 $\ddot x=-(4\pi G\rho/3)x$, 穿过地球的理想小体作简谐运动, 周期 $2\pi\sqrt{3/(4\pi G\rho)}$.

卡文迪许扭秤用已知大小与间距的球测引力. 扭丝力矩为 $-\kappa\theta$, 静平衡 $\kappa\theta=\tau_g$. 若两个小球质量 $m$, 大球质量 $M$, 距离 $d$, 力臂 $l$, 忽略其他几何修正则 $\tau_g\simeq2lGMm/d^2$. 扭转自由振荡周期 $T=2\pi\sqrt{I/\kappa}$ 可独立校准 $\kappa$. 实际需计入所有球之间的作用、距离误差与背景扰动; 引力质量与惯性质量相等是实验事实, 不是量纲恒等式.

电荷间库仑力形式 $F=k_e q_1q_2/r^2$ 同样为中心力, 但符号可吸引或排斥, 电荷不是质量. 这里只用它比较逆平方模型; 电场、电势和介质的完整内容接电磁学册.

\originalsection{有效势与圆轨道}
中心力 $\vct F=F(r)\uv r$ 对原点无力矩, $\vct L$ 守恒, 轨道限于垂直 $\vct L$ 的平面 (零角动量轨道退化为径向). $L=mr^2\dot\theta$, 面积速度 $\dot A=r^2\dot\theta/2=L/(2m)$ 恒定, 这就是开普勒第二定律的力学根源.

动能分解为径向和切向, 得
\begin{equation}
E=\tfrac12m\dot r^2+U_{\rm eff}(r),\qquad
U_{\rm eff}(r)=U(r)+\frac{L^2}{2mr^2}.
\end{equation}
最后一项称为离心势垒, 由消去角速度得到, 不是惯性系中新增的真实力. 圆轨道满足 $U'_{\rm eff}(r_0)=0$, 小径向振动频率平方为 $U''_{\rm eff}(r_0)/m$, 需二阶导数为正. 万有引力下 $v_c^2=GM/r$, $E_c=-GMm/(2r)$, 逃逸最小速度 $v_e=\sqrt{2GM/r}$, 约定无穷远速度趋于零. 逃逸能量比当前圆轨道能量高 $GMm/(2r)$, 对应速率乘 $\sqrt2$.

\originalsection{逆平方轨道与椭圆}
令 $u=1/r$, 用 $\dot\theta=Lu^2/m$ 把径向方程改写为 Binet 形式
\begin{equation}
u''+u=-\frac{mF(1/u)}{L^2u^2},
\end{equation}
撇号在此表示角度导数. 推导时 $\dot r=-(L/m)u'$, $\ddot r=-(L^2/m^2)u^2u''$, 径向加速度还减 $r\dot\theta^2=(L^2/m^2)u^3$. 引力 $F=-GMmu^2$ 给 $u''+u=GMm^2/L^2$, 所以
\begin{equation}
r=\frac{p}{1+e\cos(\theta-\theta_0)},\quad
p=\frac{L^2}{GMm^2},\quad
 e^2=1+\frac{2EL^2}{G^2M^2m^3}.
\end{equation}
最后一式可将转向点 $\dot r=0$ 的能量代入取得. $0\le e<1$ 为椭圆, $e=1$ 为抛物线, $e>1$ 为双曲线. 对椭圆 $p=a(1-e^2)$, 半短轴 $b=a\sqrt{1-e^2}$, 近日与远日半径分别 $a(1-e)$、$a(1+e)$. 吸引中心位于焦点, 不在椭圆中心.

一个周期扫过面积 $\pi ab$, 面积速度恒定, $T=2\pi mab/L$. 由 $L^2=GMm^2a(1-e^2)$ 消去 $b,L$, 得 $T^2=4\pi^2a^3/(GM)$, 即开普勒第三定律. 对两体质量可比的情形, 用相对坐标和约化质量推导, 引力参数变为 $G(M+m)$; 不能始终把主星当静止.
\begin{center}
\begin{tikzpicture}[>=Latex]
\draw (0,0) ellipse(2.2 and 1.5);\fill (.0,0) circle(1pt) node[below]{$C$};
\fill (1.61,0) circle(2pt) node[below]{$M$};\coordinate(P) at (1.1,1.30);
\fill (P) circle(2pt) node[above]{$m$};\draw[->] (1.61,0)--(P) node[midway,left]{$r$};
\draw[dashed](-2.2,0)--(2.2,0);\node at (-.8,-1.85){$a=2.2,\ b=1.5$ (示意比例)};
\end{tikzpicture}
\end{center}
图以任取比例显示焦点与半轴. 用于解释几何, 不作为某一实际行星的数据.
\originalsection{两体轨道与能量分类}
设两体位置 $\vct r_1,\vct r_2$, 质量 $m_1,m_2$, $\vct r=\vct r_1-\vct r_2$, 总质量 $M$. 分别写
\begin{equation}
m_1\ddot{\vct r}_1=-Gm_1m_2\vct r/r^3,\qquad
m_2\ddot{\vct r}_2=+Gm_1m_2\vct r/r^3.
\end{equation}
相加给质心 $\ddot{\vct R}=0$, 相减再乘约化质量 $\mu=m_1m_2/M$ 得 $\mu\ddot{\vct r}=-Gm_1m_2\vct r/r^3$. 相对运动的引力参数为 $GM$, 能量与角动量分别为 $E_{\rm rel}=\mu\dot r^2/2+L_{\rm rel}^2/(2\mu r^2)-Gm_1m_2/r$, $\vct L_{\rm rel}=\mu\vct r\times\dot{\vct r}$. 从质心定义还原
\begin{equation}
\vct r_1=\vct R+(m_2/M)\vct r,\qquad
\vct r_2=\vct R-(m_1/M)\vct r.
\end{equation}
两体相对质心的轨道是相反方向、比例 $m_2:m_1$ 的相似圆锥曲线, 椭圆周期相同. 大质量主星近似仅是 $m_2\ll m_1$ 的极限.

固定非零角动量, 写引力系数 $k=Gm_1m_2$, 有效势最低能量 $E_{\min}=-\mu k^2/(2L_{\rm rel}^2)$. $E<E_{\min}$ 不可实现; 等号为圆; $E_{\min}<E<0$ 为椭圆; $E=0$ 为抛物线逃逸; $E>0$ 为双曲线散射. $L=0$ 的径向碰撞不由这个有势垒分类完整描述.

椭圆几何也可从 $r=p/(1+e\cos\theta)$ 推出. 令焦点为原点且 $x=r\cos\theta$, 得 $r=p-ex$, 平方并配方:
\begin{equation}
\frac{(x+pe/(1-e^2))^2}{[p/(1-e^2)]^2}
+\frac{y^2}{p^2/(1-e^2)}=1.
\end{equation}
因此 $a=p/(1-e^2)$, $b=p/\sqrt{1-e^2}$, 中心在 $-ae$ 处, 焦距 $c=ae$, $b^2=a^2-c^2$. 相对轨道半长轴 $a$ 的两体周期为 $T^2=4\pi^2a^3/[G(m_1+m_2)]$, 不是某一个体半长轴直接代入固定主星公式.

\originalsection{两次脉冲的转移轨道}
要把小卫星从半径 $r_1$ 的圆轨道转到更高 $r_2>r_1$, 不是把总能量一次增加到目标值就自动得到目标圆轨道. Hohmann 转移采用与两圆相切的椭圆, 近地点 $r_1$、远地点 $r_2$, 半长轴 $a_t=(r_1+r_2)/2$. 令 $\kappa=GM$, 固定中心体且忽略卫星质量. 从 $E=-\kappa m/(2a_t)=mv^2/2-\kappa m/r$ 得活力公式 $v^2=\kappa(2/r-1/a_t)$. 两处圆速分别为 $v_{c1}=\sqrt{\kappa/r_1},v_{c2}=\sqrt{\kappa/r_2}$, 转移椭圆两端速度为
\[v_p=\sqrt{\frac{2\kappa r_2}{r_1(r_1+r_2)}},\qquad v_a=\sqrt{\frac{2\kappa r_1}{r_2(r_1+r_2)}}.\]
第一次在低圆轨道切向加速 $\Delta v_1=v_p-v_{c1}>0$, 抵高轨道后再切向加速 $\Delta v_2=v_{c2}-v_a>0$, 否则会沿椭圆返回. 两脉冲之间飞行半周期 $t_t=\pi\sqrt{a_t^3/\kappa}$. 原教材示例取 $r_1=2R_E,r_2=4R_E$, 因而 $a_t=3R_E$; 把该值代入即得到两次速度改变量和转移时间. 瞬时脉冲、同平面、两体近似是这些式的条件, 不是任意推进限制下的普遍最优性证明.

对一般二体相对轨道, 约化质量 $\mu=m_1m_2/(m_1+m_2)$, $k=Gm_1m_2$, 椭圆半长轴 $a$、偏心率 $e$ 时
\[E_{rel}=-\frac{k}{2a},\qquad L_{rel}^2=\mu k a(1-e^2).\]
在近日点 $r_p=a(1-e)$、远日点 $r_a=a(1+e)$, 径向速率为零, 角动量 $L_{rel}=\mu r v$ 直接给 $v_p=L_{rel}/[\mu a(1-e)]$, $v_a=L_{rel}/[\mu a(1+e)]$. 这与前面的相对运动重建到两体质心轨道相配合, 不能把相对速率当成任一物体相对质心速率.
